%% kapa: exam=2 round=31 year=2020 period=1 subject=appraisal_practice subject_ko=감정평가실무 seq=4 number=4 points=10
%% source: Q-Net 감정평가사 기출문제 (cst00309) artlSeq=5207673 file=raw/1008/5207673/1교시_감정평가실무.pdf
%% conversion: kordoc 4.24.51 -> tools/gen2.py (auto); review=proofread-gpt-astra; flags=figure
\begin{longquestion}{【문제 4】}
감정평가사甲은 A시 B구에서 시행하는 도시계획시설 도로 사업에 편입되는 주식회사K의 영업보상(휴업)에 대한 협의를 위한 감정평가를 의뢰받았다. 관련 법규 및 이론을 참작하고 제시된 자료를 활용하여 영업손실 보상액을 감정평가 하시오. (10점) \par
< 자료 1 > 사업의 개요 \par
1. 사업시행지: A시 B구 C동 5-19번지 일원 \par
2. 사업의 종류: 도시계획시설(도로)사업(중로5-24호선) 개설공사 \par
3. 사업시행자: B구청장 \par
4. 사업인정고시일: 2020. 01. 24 \par
5. 가격시점: 2020. 09. 19 \par
< 자료 2 > 사업 토지 및 영업장 개황 \par
\qfigure[0.65\linewidth]{figures/04-1.png}
< 자료 3 > 관련자료 \par
1. 주식회사K는 2012년 6월경 개업하였으며, 최근 3년의 월평균 영업이익은 다음과 같음 \par
{\small\centering\renewcommand{\arraystretch}{1.25}
\resizebox{\ifdim\width>\linewidth\linewidth\else\width\fi}{!}{
\begin{tabular}{|c|c|c|c|}
\hline
\makecell{구분} & \makecell{2017년} & \makecell{2018년} & \makecell{2019년} \\
\hline
\makecell{월평균 영업이익} & \makecell{3,650,000원} & \makecell{3,950,000원} & \makecell{4,250,000원} \\
\hline
\end{tabular}}\par}

2. 발전기실 \par
\medskip\noindent (1) 구조: 벽돌조 슬래브지붕 \par
\medskip\noindent (2) 연면적: 9m\textsuperscript{2} \par
\medskip\noindent (3) 사용승인일: 2012. 05. 24 \par
< 자료 4 > 기타 참고사항 \par
1. 주식회사K는 「공익사업을 위한 토지 등의 취득 및 보상에 관한 법률(이하 “토지보상법”)」 시행규칙 제45조의 영업손실의 보상대상 요건을 갖추었음 \par
2. 주식회사K는 발전기실 철거 후 재설치까지 공장가동이 불가능한 상태이며, 발전기실을 동일규모로 새로이 건축하고, 내부에 소재하는 발전기 및 그 부대설비를 이전 재설치 하는데 1개월이 소요될 예정임 \par
3. 벽돌조 슬래브지붕의 발전기실를 신축하는데 통상 1,300,000원/m\textsuperscript{2}가 소요됨 \par
4. 발전기실 내 발전기 및 그 부대설비를 이전 재설치 하는데 3,500,000원이 소요되며, 시운전 비용 500,000원이 추가 소요됨 \par
5. 도시근로자가구 월평균 가계지출비(3인가구): 4,233,829원 \par
6. 영업규모 축소에 따른 영업용 고정자산, 원재료, 제품 및 상품 등의 매각에 따른 손실은 없음 \par
7. 발전기실은 지장물 조서 목록에 별도로 조사되어 있음 \par
8. 본 영업장의 이전에 따른 휴업보상액(「토지보상법」 제47조 제1항)은 25,000,000원 임 \par
\end{longquestion}
